% id: 150-07
% book: 베이직쎈 대수 · 08 삼각함수의 활용
% section: 실전 감각 UP
% figure: tikz:fig-150-07
% answer: 
% answer_source: 
% variant_level: 0 (원본 전사)
% 이 파일은 build.mjs 가 items.json 에서 만든다. 고칠 때는 items.json 을 고친다.
\smash{\raisebox{1.5mm}{\dmkichul{베이직쎈 대수 150쪽 07번}}}
\noindent\begin{minipage}[t]{0.58\linewidth}\vspace{0pt}\raggedright 오른쪽 그림과 같이 원에 내접하는 사각형~$\pt{ABCD}$에서 $\seg{AB}=\seg{AD}=6$, $\seg{BC}=10$, $\seg{CD}=4$, $B=60^\circ$일 때, 사각형~$\pt{ABCD}$의 넓이는?\end{minipage}\hfill\begin{minipage}[t]{0.4\linewidth}\vspace{0pt}\centering \resizebox{\ifdim\width>\linewidth\linewidth\else\width\fi}{!}{% 150-07(150쪽) — 원에 내접하는 사각형 ABCD · $\seg{AB}=\seg{AD}=6$, $\seg{BC}=10$, $\seg{CD}=4$, $B=60^\circ$.
% 스캔 화질이 낮아 TikZ 로 다시 그림. 원문과 같이 B 를 왼쪽, C 를 오른쪽, A 를 왼쪽 위, D 를 오른쪽 위에 두고 안쪽을 연하게 칠했다.
% 라벨은 원문에 있는 A·B·C·D·$6$·$6$·$10$·$4$·$60^\circ$ 만 둔다 — 사각형의 넓이(답)는 적지 않는다.
\begin{tikzpicture}[line width=0.5pt, scale=1.5, >=stealth]
  \coordinate (O) at (0,0);
  \coordinate (A) at (115:1);
  \coordinate (B) at (186:1);
  \coordinate (C) at (355:1);
  \coordinate (D) at (47:1);
  \fill[red!12] (A) -- (B) -- (C) -- (D) -- cycle;
  \draw[line width=0.6pt] (O) circle (1);
  \draw[line width=0.55pt] (A) -- (B) -- (C) -- (D) -- cycle;
  \draw[line width=0.4pt] ([shift=(0.5:0.3)]B) arc (0.5:60.5:0.3);
  \node[font=\footnotesize] at ([shift=(30:0.56)]B) {$60^\circ$};
  \draw[densely dotted, line width=0.35pt, <->] ([shift=(150:0.13)]B) to[bend left=12]
    node[midway, above left=-3.5pt and -2.5pt, font=\footnotesize] {$6$} ([shift=(150:0.13)]A);
  \draw[densely dotted, line width=0.35pt, <->] ([shift=(-99:0.13)]A) to[bend right=12]
    node[midway, below=-2.5pt, font=\footnotesize] {$6$} ([shift=(-99:0.13)]D);
  \draw[densely dotted, line width=0.35pt, <->] ([shift=(-159:0.13)]D) to[bend right=12]
    node[midway, left=-1.5pt, font=\footnotesize] {$4$} ([shift=(-159:0.13)]C);
  \draw[densely dotted, line width=0.35pt, <->] ([shift=(-89:0.13)]B) to[bend right=10]
    node[midway, below=-2.5pt, font=\footnotesize] {$10$} ([shift=(-89:0.13)]C);
  \node[above left=-3.5pt and -3pt, font=\small] at (A) {$\pt{A}$};
  \node[left=1pt, font=\small] at (B) {$\pt{B}$};
  \node[right=1pt, font=\small] at (C) {$\pt{C}$};
  \node[above right=-3pt and -2.5pt, font=\small] at (D) {$\pt{D}$};
\end{tikzpicture}
}\end{minipage}\par
\dmchoicesauto{\rule[-3ex]{0pt}{8ex}}{$28$}{$20\sqrt{2}$}{$21\sqrt{2}$}{$20\sqrt{3}$}{$21\sqrt{3}$}
